\section{The cubic second moment and its modular comparison}\label{sec:cubic}
Throughout this section the character has order three. Put
\[
 \mathcal G_2=\Sym^2\mathcal V\otimes\mathcal L_\rho,
 \qquad W_2=H^1_!(U_{\overline K},\mathcal G_2).
\]
The preceding geometric construction gives $\dim_EW_2=2$, unramifiedness
outside the primes above $2,3$, and, for every other prime ideal
$\mathfrak p$ with $q=N\mathfrak p$,
\begin{equation}\label{eq:cubic-trace}
 \Tr(\Frob_{\mathfrak p}\mid W_2)=-M_{3,2}(q)-q-1.
\end{equation}
Here Frobenius is geometric, $\tau=E(-1)$ has value $q$, and the Kummer
normalization is the actual torsor $z^3=\lambda$. We first obtain the
remaining determinant constraint geometrically. The modular comparison
then uses the separately recorded finite arithmetic certificate.

\subsection{Arithmetic inversion and the determinant}
Let $i(\lambda)=\lambda^{-1}$. On the etale cover $u^2=\lambda$, the map
\[
 \phi_u:\mathcal E_\lambda\longrightarrow\mathcal E_{1/\lambda},
 \qquad (x,y)\longmapsto (x/\lambda,y/u^3)
\]
is an elliptic-curve isomorphism: substitution gives
$(y/u^3)^2=(x/\lambda)(x/\lambda-1)(x/\lambda-1/\lambda)$.
Replacing $u$ by $-u$ composes with $[-1]$, whose action on
$\Sym^2\mathcal V$ is trivial. Thus the induced isomorphism
$i^*\Sym^2\mathcal V\simeq\Sym^2\mathcal V$ descends arithmetically.
The map $z\mapsto z^{-1}$ on the cubic cover inverts the deck group,
so $i^*\mathcal L_\rho\simeq\mathcal L_\rho^\vee$, $E$-linearly over $K$.
This step does not conjugate the coefficient field.

The elliptic cup product induces the nondegenerate symmetric pairing
\[
 (v_1v_2,w_1w_2)\longmapsto
 \langle v_1,w_1\rangle\langle v_2,w_2\rangle+
 \langle v_1,w_2\rangle\langle v_2,w_1\rangle
 \quad\text{with values in }E(-2).
\]
Consequently $i^*\mathcal G_2\simeq\mathcal G_2^\vee(-2)$.
Poincare duality \cite[Theorem 24.1]{Milne} induces a perfect pairing
on parabolic images: the kernels of the maps $H_c^1\to H^1$ are the
annihilators of the corresponding dual images. Hence
\[
 H^1_!(U,\mathcal G_2^\vee)\simeq W_2^\vee(-1),
 \qquad W_2\simeq W_2^\vee(-3).
\]
The inversion pullback commutes with $G_K$ and with the compact-support
map, so these are arithmetic isomorphisms. Taking determinants gives
\begin{equation}\label{eq:cubic-det-sign}
 \det W_2=\tau^3\varepsilon,\qquad \varepsilon^2=1,
\end{equation}
where $\varepsilon$ is continuous and unramified outside $2,3$.
In usual cyclotomic notation this is
$\chi_{\rm cyc}^{-3}\varepsilon$. The assertion
$\varepsilon=1$ will use the finite comparison below.

\subsection{The modular representation and residual conditions}
Let $f=f_{36}$ be the weight-four normalized newform of level $36$
and trivial nebentypus, labelled $36.4.a.a$. Its eta expression is
\cite[Table 1]{LTYZ}
\begin{equation}\label{eq:f36-eta}
 f(z)=\frac{\eta(6z)^{14}}{\eta(2z)^3\eta(18z)^3}
       -3\eta(2z)^3\eta(6z)^2\eta(18z)^3.
\end{equation}
Use the cohomological representation $R_f$ over $E$: it is the dual of
the conventional Deligne representation normalized at arithmetic
Frobenius; see the modular-form motive and its etale realization in
\cite[Theorem~1.2.4(i)]{SchollMotives}.
Thus geometric Frobenius at $p\ge5$ on $R_f$ has trace
$a_f(p)$ and determinant $p^3$. Both $W_2$ and $R_f|_{G_K}$ are
unramified outside primes above $2,3$.

Their compact images admit stable $\mathcal O$-lattices, with
$\mathcal O/(2)=\F_4$. Full rational $2$-torsion on every smooth
Legendre fibre over $\F_q$ makes $a_q(\lambda)$ even. Therefore
$M_{3,2}(q)\in4\mathcal O$, and
\eqref{eq:cubic-trace}--\eqref{eq:cubic-det-sign} give
\[
 \Tr W_2\equiv0\pmod2,\qquad \det W_2\equiv1\pmod2.
\]
Initially these congruences hold on good Frobenius classes;
Chebotarev density and continuity extend them to $G_K$.

The carried certificate records, and the independent calculation described
below rechecks,
$f\equiv\eta(3z)^8\pmod2$ through the Sturm bound
\[
 \frac4{12}[\operatorname{SL}_2(\Z):\Gamma_0(36)]
 =\frac4{12}\,72=24.
\]
Here $\eta(3z)^8\in S_4(\Gamma_0(9))\subset S_4(\Gamma_0(36))$
by \cite[Table~1]{LTYZ}. Both sides have integral Fourier coefficients.
Sturm's theorem \cite[Theorem~1, p.~276]{Sturm} makes this a full
Fourier-series congruence,
subject to that recorded finite check. With $Q_z=e^{2\pi iz}$,
\[
 \eta(3z)^8\equiv\eta(24z)
 =\sum_{m\ge1}\chi_{12}(m)Q_z^{m^2}\pmod2.
\]
The congruence is the characteristic-two product identity; the equality
is Euler's pentagonal-number identity, with
$\chi_{12}(m)=(12/m)$ the Kronecker symbol. Prime-index coefficients of $f$ are therefore
even. At a split prime of $K$, the restricted modular trace is $a_f(p)$;
at an inert prime it is $a_f(p)^2-2p^3$, again even. Its determinant
is $q^3$. Density gives the required global congruences for
$R_f|_{G_K}$ as well. The residue field $\F_4$ is allowed in the
comparison criterion used below.

\subsection{The carried finite certificate and its consequence}
Let $K_S$ be the compositum of the quadratic extensions of $K$
unramified outside the primes above $2,3$. The Eisenstein integer ring
has class number one. Removing square factors from an allowed
squareclass therefore leaves an $S$-unit. Its unit squareclass is
generated by $-1$, and the primes above $2$ and $3$ can be generated
by $2$ and $\sqrt{-3}$. Their valuations, together with the nonsquare
unit $-1$, prove independence. Consequently
\[
 K_S=K\bigl(\sqrt{-1},\sqrt2,\sqrt{\sqrt{-3}}\bigr),
 \qquad\Gal(K_S/K)\simeq\F_2^3.
\]
The seven nonzero vectors constitute a non-cubic set: a homogeneous
cubic polynomial function vanishing on them also vanishes at zero,
and hence vanishes on the whole vector space.

Table~\ref{tab:cubic-certificate} reproduces the cubic rows of the
carried exact certificate. A prime ideal is specified by $p$ and the
image $r$ of $\sqrt{-3}$ modulo it. A bit is $0$ for a square and $1$
for a nonsquare; the displayed order is
$(\sqrt{-3},2,-1)$, the reverse of the basis above.
\begin{table}[htbp]
\centering
\begin{tabular}{rrrrr}
\hline
$p$ & $r$ & bits & $t=a_f(p)$ & $d$ \\
\hline
7  & 2  & 001 & 8    & 343 \\
7  & 5  & 101 & 8    & 343 \\
13 & 6  & 110 & $-10$  & 2197 \\
19 & 4  & 011 & $-100$ & 6859 \\
19 & 15 & 111 & $-100$ & 6859 \\
37 & 16 & 010 & $-226$ & 50653 \\
73 & 17 & 100 & 26   & 389017 \\
\hline
\end{tabular}
\caption{Previously recorded exact cubic comparisons, with $d=p^3$.
The seven rows and Sturm comparison have also been recomputed by a new
standalone exact implementation.}
\label{tab:cubic-certificate}
\end{table}

The certificate consists of the supplied
\texttt{V2\_01\_exact\_livne.py}, its exact arithmetic helpers,
\texttt{V2\_exact\_livne.txt}, and
\texttt{exact\_livne\_certificate.json}. Static inspection shows direct
finite-field character sums and integer eta-series arithmetic for the
cubic comparisons, with no FFT rounding. The certificate uses the
norm-lifted character over $\F_{p^2}$. Thus, by the geometric trace
formula,
\begin{align*}
 t&=-M_{3,2}(p)-p-1,\\
 u&=-M_{3,2}(p^2)-p^2-1=\Tr(\Frob_{\mathfrak p}^{2}\mid W_2),\\
 d&=\frac{t^2-u}{2}.
\end{align*}
The inversion identity for even Legendre moments makes the two
primitive cubic-character choices equivalent here. The recorded
values satisfy $t=a_f(p)$ and $d=p^3$ at every listed prime; the
recorded Sturm comparison reaches degree $24$.
\emph{The archived programs were not rerun.} A new standalone implementation,
\texttt{independent\_cubic\_exact.py}, independently recomputes the
displayed cubic comparisons using direct finite-field character sums and
integer eta-series arithmetic. It uses only the Python standard library;
the mathematical calculations use arbitrary-precision integers, without
FFT or floating-point arithmetic. Over $\F_{p^2}=\F_p[s]/(s^2-d)$ it uses
$\chi_{p^2}(a+bs)=\chi_p(a^2-db^2)$ and the norm-lifted cubic character.
The new log records agreement at the five rational primes underlying the
seven rows, all seven square-class vectors, and the congruence through
degree $24$. Historical expected values are read only after the new
calculation, for comparison. Checks include every rational-fibre quadratic
base-change identity and a second direct multiplication calculation over
$\F_{49}$. Source hashes and the new log are supplied in the release
supplement. This rechecks the bounded arithmetic; it does not re-prove the
geometric hypotheses, newform identification, or external comparison
theorems. The theorem below retains its explicit finite-check dependencies.

\begin{theorem}[Cubic modular comparison using the carried certificate]
\label{thm:cubic-modular}
Assuming the finite exact comparisons and the Sturm comparison just
specified, the representations have isomorphic semisimplifications:
\[
 W_2^{\rm ss}\simeq(R_f|_{G_K})^{\rm ss},
 \qquad \det W_2=\tau^3.
\]
For every good prime ideal $\mathfrak p$ of $K$,
\[
 M_{3,2}(N\mathfrak p)
 =-\Tr(\Frob_{\mathfrak p}\mid R_f|_{G_K})-N\mathfrak p-1.
\]
In particular, for $p\ge5$ and $p\equiv1\pmod3$,
$M_{3,2}(p)=-a_f(p)-p-1$; for $p\equiv2\pmod3$,
\[
 M_{3,2}(p^2)=-a_f(p)^2+2p^3-p^2-1.
\]
\end{theorem}
\begin{proof}
Livne's criterion in \cite[Theorem 4.1]{DGP} applies to continuous
rank-two representations over a finite extension of $\Q_2$.
The common ramification set, vanishing residual traces, and equal
residual determinants were established above. The seven rows give
exact trace and determinant agreement on a non-cubic set in
$\Gal(K_S/K)$. The criterion gives the asserted semisimplification.
If it is formulated using the opposite Frobenius convention, invert
both matrices: $\Tr(A^{-1})=\Tr(A)/\det A$, and the quadratic test
classes are unchanged under inversion.

Alternatively to taking the determinant of the comparison,
\eqref{eq:cubic-det-sign} and the table give
$\varepsilon(\Frob_{\mathfrak p})=d/p^3=1$ on all seven nonzero
classes. Since $\varepsilon$ factors through $\Gal(K_S/K)$, it is
trivial. Finally apply \eqref{eq:cubic-trace}; restriction at an inert
rational prime replaces Frobenius by its square.
\end{proof}

This section proves the stated cubic second-moment comparison with its
explicit certificate dependency. It supplies no geometric proof for
characters of orders four or six, and no bad-prime Euler-factor or
rank-four modularity assertion.
